\clearpage
\VerdierPhysicalPageAnchor{0190}%
\begingroup\small\setlength{\parskip}{0pt}\setlength{\emergencystretch}{1.5em}
\setlength{\abovedisplayskip}{4pt plus 2pt minus 2pt}%
\setlength{\belowdisplayskip}{4pt plus 2pt minus 2pt}%

\VerdierSemanticAnchor{ch3d-2-1-5}
\noindent\textbf{Remarque 2.1.5.}\quad
Plaçons nous sous les hypothèses du lemme (2.1.4) et supposons de plus
que $X,Y\in\operatorname{Ob}\operatorname{comp}^{-}(\mathcal A)$.
L'argument présenté plus haut montre alors qu'on peut, en changeant au
besoin $X$ par un homotopisme, supposer que $m$ est un épimorphisme. La
construction de la suite $X_n$ présentée ci-dessus, montre alors que le
système inductif $X_n$ est stationnaire (i.e. $X_n=Y$ et
$u_n=\operatorname{identité}$ pour $n$ petit) et que pour tout $n$,
$X_n\in\operatorname{Ob}\operatorname{comp}^{-}(\mathcal A)$.

\VerdierSemanticAnchor{ch3d-2-2}
\subsection{Étude de certains ensembles d'objets.}

\VerdierSemanticAnchor{ch3d-2-2-1}
\noindent\textbf{2.2.1.}\quad
Soit $M$ un ensemble d'objets de $\mathcal A$. Nous considérerons les
différentes propriétés suivantes :

\VerdierSemanticAnchor{ch3d-2-2-1-e1}
\noindent$(E_1)$ Les objets nuls appartiennent à $M$. La somme directe de
deux objets appartenant à $M$ est un objet appartenant à $M$.

\VerdierSemanticAnchor{ch3d-2-2-1-e2}
\noindent$(E_2)$ Pour tout épimorphisme
$X\xrightarrow{u}L$, où $X$ est un objet de $\mathcal A$ et $L$ un élément
de $M$, il existe un épimorphisme $v\colon L'\longrightarrow L$,
$L'\in M$, qui se factorise par $u$.

\VerdierSemanticAnchor{ch3d-2-2-1-e3}
\noindent$(E_3)$ Pour tout entier $n$ et toute suite exacte :
\[
0\longrightarrow X\longrightarrow L_n\longrightarrow L_{n-1}
\longrightarrow\cdots\longrightarrow L_0\longrightarrow0,
\qquad L_i\in M,\quad0\leq i\leq n,
\]
il existe un épimorphisme $L'\longrightarrow X$, où $L'\in M$.

\VerdierSemanticAnchor{ch3d-2-2-1-e4}
\noindent$(E_4)$ Pour tout épimorphisme
$u\colon L\longrightarrow L'$, $L,L'\in M$, le noyau de $u$ est un
élément de $M$.

\VerdierSemanticAnchor{ch3d-2-2-1-e5}
\noindent$(E_5)$ Pour toute extension
$0\longrightarrow L\longrightarrow Z\longrightarrow L'\longrightarrow0$,
$L,L'\in M$, l'objet $Z$ est un élément de $M$ ; $0$ est un élément de
$M$.

\VerdierSemanticAnchor{ch3d-2-2-1-e6}
\noindent$(E_6)$ L'ensemble $M$ est exact, i.e.\ chaque fois que deux des
objets d'une suite exacte courte de $\mathcal A$ sont des éléments de $M$,
le troisième l'est aussi ; $0$ est un élément de $M$.

\VerdierSemanticAnchor{ch3d-2-2-1-implications}
On notera qu'on a les implications $(E_6)\Rightarrow(E_4)$ et $(E_5)$,
$(E_4)\Rightarrow(E_3)$ et $(E_5)\Rightarrow(E_1)$. Une suite exacte
illimitée :

\VerdierSemanticAnchor{ch3d-2-2-1-infinite-presentation}
\[
\cdots\longrightarrow L_n\longrightarrow L_{n-1}
\longrightarrow\cdots\longrightarrow L_0\longrightarrow X
\longrightarrow0,
\qquad L_i\in M,\quad i\in\mathbb N,
\]
est appelée une $\infty$-$M$-présentation de l'objet $X$. Soit $M$ un
ensemble d'objets possédant les propriétés $(E_1)$, $(E_2)$ et $(E_3)$.
Désignons par $M^{\wedge}$ l'ensemble d'objets suivants : un objet $X$
appartient à $M^{\wedge}$ s'il possède une $\infty$-$M$-présentation et si,
pour tout entier $n$ et toute suite exacte :

\VerdierSemanticAnchor{ch3d-2-2-1-m-wedge}
\[
0\longrightarrow X'\longrightarrow L_n\longrightarrow L_{n-1}
\longrightarrow\cdots\longrightarrow L_0\longrightarrow X
\longrightarrow0,
\qquad L_i\in M,\quad0\leq i\leq n,
\]
l'objet $X'$ possède une $\infty$-$M$-présentation.
\endgroup
